\section{The collapsing end and its action jump}\label{sec:jump}

We compute the difference of the two stationary actions at the same
parameter $p=(b,x)$, hence at the same compact observation value
$\mathsf P(p)$ in the notation of \Cref{sec:arrays}.
The two actions are evaluated on solutions of the
complexified stationary equation.  In particular, the translator that
appears below is a limit of exterior complex boundary problems; no
translating cap in the mean curvature flow is assumed.

Write $B=|a|^{-1/2}$, and fix the contour constants before letting
$B\longrightarrow\infty$.  An adverse end has opening coefficient $\beta$
and
\begin{equation}\label{jump:pole}
 Z_0=(-\beta)^{-1/2},\qquad |Z_0|\asymp B,
 \qquad r^2/2=1-z^2/Z_0^2.
\end{equation}
The branch of $Z_0$ is near the positive real axis.  A sufficiently small
fixed change of the parameter's evaluation angle is permitted.  All
constants below are uniform under that change and on the finite set of
adverse ends.  Squares, inner products and square roots are complex
bilinear, with branches continued from the positive real area element.

The leading pole square in \eqref{arr:A0} is $(-c_ja)^{-1}$,
whereas \eqref{jump:pole} uses the actual first-jet square
$-1/\beta_j(p(s))$ at the common parameters.  Their ratio is
$c_ja/\beta_j(p(s))=1+o(1)$ on the fixed admissible ray and finite
adverse end list.  This relative comparison preserves strict margins
in rates proportional to $B^2$; it does not identify their Gaussian
exponentials up to a relative $1+o(1)$ factor.  The mean-matching radius
$Z$ and shifted exterior radius $Z_*$ are distinct again, defined in
\eqref{jump:matching-Z} and \eqref{jump:scaling} below.

We use the following precise interfaces from the preceding sections.
The preliminary normalized solutions of
\cref{arr:preliminary-solves}, for the holomorphic comparison-data
families used here, have discrepancy $O(B^{-1.8})$ in the strip
and trace norms on fixed scaled bands, and are holomorphic in the finite
parameters on balls of fixed power radius.  The contours of
\cref{arr:contours} have strict diffusion margins and allow shortening
and the height turns described below.  The estimates in
\cref{lin:strip,lin:shielding,lin:global} provide, respectively, local
strip and boundary estimates, exponential suppression of far boundary
data on a fixed overlap, and the normalized global inverse with growing
weights retained until the final estimate.  Finally,
\cref{gauge:mixed,gauge:height} apply to the full commuted nonlinear
operators in the prescribed mixed and height gauges.  We spell out the
weights and the small quantities needed in each application; an
unweighted polynomial inverse is used only for an exponentially small
residual.

We first synchronize the complete nonadverse boundary problems at the
full stop $\Re z^2=DB^2$.  On the chosen evaluation ray every nonzero
nonadverse leading pole is on the favorable side, so in the normalized
coordinate of \cref{arr:contours} its direction $p_*$ satisfies
$\Re(1/p_*)<0$, with a fixed margin on the nearby rays in use.  The
favorable part of that proof first deforms each profile to an admissible
ray.  A ray $u=h(1-im)$ is admissible there when
$\Re((1-im)/p_*)<0$; this inequality is affine in $m$.  Interpolating
its slope to $m=0$ therefore keeps the strict inequality and the same
horizontal stop.  The initial real collars are held fixed, and the
strict tangent conditions permit the stated smoothing at their joins.
Smaller-order openings already use the real path.  First, on each
original fixed path, interpolate any permitted remainder datum linearly
to zero in its normalized natural trace ball
\eqref{arr:allowed-traces}, which is convex.  Then deform the path while
prescribing the outgoing derivative of its varying comparison profile.
Both arrays consequently reach the same nonadverse path and the same
comparison-profile outgoing derivative on each end.  All these
normalized data and path derivatives have fixed polynomial bounds.

This synchronization is an actual change of the normalized actions.
Write $G_\pm^{\rm syn}$ for the resulting actions.  Preliminary
uniqueness supplies the normalized homotopies, and
\cref{arr:first-variation} gives
\begin{equation}\label{jump:nonadverse-sync}
 |G_\pm-G_\pm^{\rm syn}|
       \le C_D B^M e^{-DB^2/5}.
\end{equation}
The estimate retains boundary-height variations caused by the data
change, even at a fixed exit.  Increase $D$, before choosing $B$, so
that $D/5>\max_{j\ \mathrm{adverse}}(4|c_j|)^{-1}$ with a fixed
margin.  The pole comparison and
\eqref{jump:matching-Z}--\eqref{jump:scaling} below give
$Z_{*,j}^2/Z_{j,\mathrm{lead}}^2=1+o(1)$, so the right
side of \eqref{jump:nonadverse-sync} is then negligible relative to
each absolute adverse action scale.  Below the starting arrays are
these synchronized ones.  Their nonadverse paths, exit sections, and
outgoing derivative prescriptions remain fixed and identical through
all subsequent shortening and terminal-data homotopies.  Their boundary
heights need not agree, so the corresponding far boundary momenta will
still be retained in the action estimate.

\subsection{Prescribed height charts and the matched phase}

First shorten the adverse contours while their exits satisfy
$\Re(z^2-Z_0^2)\geq cB^2$.  They can then follow the shrinking circular
profile to a small fixed radius and turn above or below zero in $r^2$,
ending near $\arg r^2=\pm2\pi/3$.  On a turn require
\begin{equation}\label{jump:turn}
 \Re\,d\log r<0,\qquad
 \Re\bigl(Z^2d(r^2)\bigr)<0.
\end{equation}
The inequalities have fixed margins after parametrization.  They say
that the angular diffusion advances and that the real fast exponent
increases.  One obtains them by starting with a radial decrease,
turning on a slightly contracting spiral, and stopping before phase
$\pi$ in $r^2$.  The homotopies from the original contours are those of
\cref{arr:contours}; at a real adverse pole their defining inequality
is affine in $Y^2$ and $(Y^2)'$ for $z^2/Z_0^2=h+iY$.  Strictness
persists under the fixed small phase change.  Thus this deformation
costs, by \cref{arr:first-variation}, at most a fixed power of $B$ times
$\exp(-\Re Z_0^2/4-cB^2)$ in action.

Choose a shared real incoming section $r=r_0$, with $r_0>0$ small and
fixed.  On fixed overlaps use one prescribed transverse parametrization
\[
 X(t,\theta)=
 \bigl(B(\zeta_0(t)+c_1(t)u),
       (R_0(t)+c_2(t)u)e_r(\theta)\bigr).
\]
The reference determinant $c_2\zeta_0'-c_1R_0'$ stays away from zero;
the terminal gauge is pure relative height.  Coefficients and cutoffs
are prescribed on the real parameter interval.  This is a single
quasilinear graph problem, and does not require inversion of an
unknown complex spatial coordinate.

Fix a real interval $0\le t\le T_{\rm pre}$ beyond $r_0$, with
$t=\log(r_0/r)$ and $T_{\rm pre}>0$ fixed, and reserve inside it a
joining band separated by fixed positive gaps from both endpoints.
Use a common preliminary solution stopped at $T_{\rm pre}$, with the
frozen nonadverse paths and data just specified, at exactly the given
parameter $p$.  All these choices precede the phase match.  Define $Z$ by
\begin{equation}\label{jump:matching-Z}
 \frac1{2\pi}\int_{\mathbb S^1}z(r_0,\theta)\,d\theta=Z\phi(r_0),
 \qquad \phi(r)=\sqrt{1-r^2/2}.
\end{equation}
Then $Z/Z_0=1+O(B^{-1.8})$.  The relative-height discrepancy at this
section has zero mean and strip and trace size $O(B^{-1.8})$.
Define $d_0$ and $Z_*$ by
\begin{align}
 (r^{-1}-r/2)d_0'+d_0/2
   &=-\phi''/(\phi')^2,\qquad d_0(r_0)=0,\label{jump:d0}\\
 d_0(r)&=2\log r+d_*+O(r^2(1+|\log r|)),
 \label{jump:dstar}\\
 Z_*&=Z+Z^{-1}\bigl(d_*-2\log(Z/2)\bigr),\label{jump:scaling}\\
 y&=Z_*r/2,\qquad h=(Z_*-z)Z_*/2.\notag
\end{align}
Logarithms are continued from positive real numbers.  There is no
existence issue in \eqref{jump:d0}: direct differentiation gives
\begin{equation}\label{jump:d0-integral}
 \left(\frac{d_0}{\phi}\right)'=\frac{2}{r\phi^4},
 \qquad
 d_0(r)=\phi(r)\int_{r_0}^{r}\frac{2\,ds}{s\phi(s)^4}.
\end{equation}
Expanding the integrand at zero proves \eqref{jump:dstar}, with $d_*$
real for the fixed real matching section.  It also controls every fixed
number of logarithmic derivatives on the chosen branches.

The leading comparison profile on the height chart is
\begin{equation}\label{jump:height-profile}
 g_c(r)=\phi(r)+Z^{-2}d_0(r),\qquad g=z/Z.
\end{equation}
Its derivative is prescribed at the exit.  Changing to that prescription
on a fixed scaled bypass is allowed by the fast boundary lift below.

For a radial graph at this finite scale, put $l=h_y$.  Its radial
area form, continued from the positive real one, is
\[
 2\pi r\,e^{-(r^2+z^2)/4}\sqrt{1+z_r^2}\,dr.
\]
Substituting $r=2y/Z_*$ and $z=Z_*-2h/Z_*$ gives exactly
\begin{equation}\label{jump:physical-action}
 \frac{8\pi}{Z_*^2}e^{-Z_*^2/4}
       ye^h\exp\left(-\frac{h^2+y^2}{Z_*^2}\right)
                                      \sqrt{1+l^2}\,dy.
\end{equation}
Thus one adverse end carries the exact action factor
$Z_*^{-2}e^{-Z_*^2/4}$.  A relative $1+o(1)$ approximation of $Z_*^2$
does not determine this exponential to relative $1+o(1)$ accuracy.
Finite relative end weights require $Z_*^2$ through an additive $o(1)$
error.  The absolute expansion in \Cref{phase:expansion} and the
full-principal-part comparison in \Cref{phase:no-arc} supply that
information after the local jump has been computed.

\begin{proposition}[Universal action jump]\label{jump:jump}
Let $G_+$ and $G_-$ be the two array actions at a common parameter
$p$ in the inner parameter ball, and hence at common compact observation
values, for a nonempty finite adverse end list with the same ordering of
the bypasses on every end.
Starting from these arrays, make the nonadverse synchronization
\eqref{jump:nonadverse-sync} and use the single normalized preliminary
solution stopped at $T_{\rm pre}$ above.  Match $Z_{*,j}$ from that
solution by \eqref{jump:matching-Z}--\eqref{jump:scaling}.  The adverse
continuations use the canonical profile $g_c$ and outgoing remainder
derivatives
\[
 \psi=\left.\partial_s(g-g_c)\right|_{\mathrm{exit}}
\]
in the natural Neumann trace class \eqref{jump:permitted-data}, in the
prescribed real path parameter at each exit.  Let $Y_0$ be the fixed large
exit cutoff for this class.  The continuation at the first fixed radius
interpolates the preliminary circular derivative and that of $g_c$ as in
Proposition~\ref{jump:compatibility}.
The longer shrinking branches have $\psi=0$ at their own exits, and the
shorter interpolation uses their actual remainder traces at the common
shorter section.
Proposition~\ref{jump:compatibility} proves membership in the relevant
uniqueness balls and exact equality of the interpolation endpoints with
the longer restrictions.

For these actions there is a universal constant $\mathfrak c\ne0$ such
that
\begin{equation}\label{jump:action-jump}
 G_+-G_-=\mathfrak c
       \sum_{j\ \mathrm{adverse}}
       Z_{*,j}^{-2}e^{-Z_{*,j}^2/4}
       +o\left(\sum_{j\ \mathrm{adverse}}
                    |Z_{*,j}|^{-2}e^{-\Re Z_{*,j}^2/4}\right).
\end{equation}
With $c_0$ defined by the upper-minus-lower $y^2$ lateral labels in
Lemma~\ref{jump:stokes},
$\mathfrak c=\epsilon_{\rm lab}8\pi c_0$ for one common
$\epsilon_{\rm lab}\in\{1,-1\}$ determined by the array labels.
Indeed an upper axial bypass in $z^2/Z_0^2$ becomes a lower
$r^2$ bypass under $r^2=2(1-z^2/Z_0^2)$; this is one global
label reversal, not an end-dependent sign.  Reversing both
array labels reverses $\mathfrak c$.  More precisely, the
limsup of the relative error as $B\to\infty$ at fixed $Y_0$
is bounded by a quantity tending to zero as $Y_0\to\infty$.
\end{proposition}

The remainder of this section constructs the stated height families,
proves the two compatibility assertions and the endpoint estimate,
and computes the common constant by the translator boundary momentum.

The exit radius will decrease to $|y|\asymp Y$, where $Y\geq Y_0$ and
$Y_0$ is a large fixed constant.  On each final bypass, $|y|\asymp Y$,
the length is $O(Y)$, and
\begin{equation}\label{jump:negative-exit}
 \Re y_{\rm e}^{2}\leq-cY^2.
\end{equation}
After $Z_*$ has been determined, put its small phase-alignment band
strictly beyond the preliminary stop $T_{\rm pre}$.  It rotates the
later incoming part so that $y$ is real there while preserving
\eqref{jump:turn}.  Every exterior path agrees with the same real
path on $[0,T_{\rm pre}]$; in particular the joining band and its
far-side gap precede this $Z_*$-dependent alignment and every terminal
turn.  Choose homothetic bypasses at intermediate scales, with bounded
derivatives in logarithmic size.  Choose every smoothing and phase-alignment
profile once in normalized log coordinates, with joins and collars of fixed
positive width; only the small phase parameters vary.  Homothetic scaling
then preserves uniform bounds for the logarithmic path tangent, its inverse,
and every required fixed derivative.

\subsection{The inward height equation and its inverse}

For completeness, the exact height equation in polar orthonormal
coordinates is
\begin{equation}\label{jump:height-equation}
 A z_{rr}+B_1\left(\frac{z_r}{r}+
                 \frac{z_{\theta\theta}}{r^2}\right)
 +2C_1\left(\frac{z_{r\theta}}r-
                 \frac{z_\theta}{r^2}\right)
 =\frac{rz_r-z}{2},
\end{equation}
where $(A,C_1;C_1,B_1)$ is
$I-Dz\otimes Dz/(1+Dz\cdot Dz)$ and
$Dz=(z_r,z_\theta/r)$.  Put $t=\log(r_0/r)$ and
$E\asymp|Br|^{-2}$.  Multiplying \eqref{jump:height-equation} by
$r^2/Z$ gives the exact expression
\begin{equation}\label{jump:log-height}
 A(g_{tt}+g_t)+B_1(g_{\theta\theta}-g_t)
 -2C_1(g_{t\theta}+g_\theta)
 +\frac{r^2}{2}(g_t+g)=0.
\end{equation}
On each bounded modified band this formula is pulled back to its
real path parameter $s$.  These bands include the phase transition
beyond $T_{\rm pre}$ making $y$ real and the terminal turn.  Put $J=t_s$.
The prescribed $J,J^{-1}$ and their fixed derivatives are bounded
on these bands, and all strip norms
there use $s$.  It is useful to record the entire drift division.
If the circular linearization in the $t$ coordinate is
\[
 (1+\widetilde b_1)v_t-a_1v_{tt}
                -a_2v_{\theta\theta}+\widetilde b_0v,
\]
then its pullback, multiplied by $J$, is
\[
 Dv_s-\frac{a_1}{J}v_{ss}
       -Ja_2v_{\theta\theta}+J\widetilde b_0v,\qquad
 D=1+\widetilde b_1+a_1\frac{J_s}{J^2}.
\]
We divide by this full nonzero coefficient $D$, including the
chain-rule term.  For the nonlinear equation and its residual, we use the
fixed circular row \eqref{gauge:fixed-row} with $g_0=g_c$ and
$\Delta_{\mathrm{ref}}=D$.  The factor $D$ is fixed at $g_c$ during graph
variation, so the derivative of that normalized map is exactly the
following circular operator even when the profile residual is nonzero:
\begin{equation}\label{jump:height-linear}
 \mathcal L_0v=v_s-a(s)v_{ss}-c(s)v_{\theta\theta}+q_0(s)v.
\end{equation}
On each long incoming piece $J=1$, so $t$ differs from $s$ only
by a fixed additive constant.

Here are the coefficient bounds, including the zeroth-order term
needed on a turn.  At circular data in
\eqref{gauge:polar-normalized}, put
\[
 P=(g_c)_t/r^2,\quad e=(Zr)^{-2},\quad
 H=(1-r^2/2)P^2-r^2e/2.
\]
The radial, angular and zeroth-order coefficients before
linearization are respectively
$a_c=e/H$, $c_c=(e+P^2)/H$, and $m_c=r^2c_c/2$.
For $|Br|\geq Y_0$, $P,H$ are bounded units,
$a_{c,P}=O(E)$, $c_{c,P}=O(E)$, and $m_{c,P}=O(r^2E)$.
More precisely, \eqref{jump:d0-integral} gives
$P=(2\phi)^{-1}+O(E)$ and hence
$H=1/4+O(r^2+E)$,
$a_c=4e(1+O(r^2+E))$, and $c_c=1+O(r^2+E)$.
Since $(g_c)_{tt}=O(r^2)$ and $g_c=O(1)$, the circular linearized
drift is
\[
 D_0=1-r^{-2}\{a_{c,P}(g_c)_{tt}+m_{c,P}g_c\}=1+O(E).
\]
The pulled drift $D_0+a_cJ_s/J^2$ is therefore a uniform unit.
In particular, after its division,
\begin{gather}
 \Re a\geq c_*E,\quad |a|\leq C_*E,\qquad
 \Re c\geq c_*,\quad |c|\leq C_*,
 \label{jump:height-coefficients}\\
 c=1+O\bigl(|r|^2(1+|\log r|)+E\bigr)
       \quad\hbox{on the long incoming pieces},\notag\\
 q_0=O(|r|^2)\quad\hbox{on the whole path}.
 \notag
\end{gather}
All fixed normalized derivatives have the same bounds.
We choose the initial phase transition as a slightly contracting
spiral, so it has the same strict conditions as the terminal turn.
On every modified band \eqref{jump:turn} says
$\Re J\geq c|J|$ and
$\Re(Z^2r^2J)\geq c|Z^2r^2J|$ with fixed relative margins.
The preceding expansions then give the
angular and radial inequalities in \eqref{jump:height-coefficients}.
The modified bands have bounded total logarithmic length and
bounded variation of $\log|r|$.  Group adjacent modifications as
one such band.  At each interface with an incoming piece reserve a
fixed $J=1$ collar, included among the incoming pieces; all ramps of
the oscillatory weight below take place on these collars.  On the long incoming pieces,
where $J=1$, $E\asymp B^{-2}|r|^{-2}$
increases toward the exit,
$\int |q_0|\,ds\leq C r_0^2$, and
$\int E\,ds\leq C Y_0^{-2}$.  The bounded modified bands preserve these
bounds up to fixed constants.  We take $r_0$ small and then $Y_0$
large.  Noncircular coefficient changes will be treated by the
projected nonlinear estimates of \cref{gauge:height}.

Here and below the strip norm $X_E^h(v;I)$ on a unit interval contains
\begin{equation}\label{jump:strip-norm}
 \|v\|_{L^2_tH^h_\theta}+\|v_t\|_{L^2_tH^h_\theta}
 +\|v_{\theta\theta}\|_{L^2_tH^h_\theta}
 +\|Ev_{tt}\|_{L^2_tH^h_\theta}
 +\|\sqrt E\,v_{t\theta}\|_{L^2_tH^h_\theta},
\end{equation}
together with the traces of $v,v_\theta,\sqrt E v_t$ at the same
fixed sufficiently high angular order.  On a strip the values of $E$
are comparable.  Source norms are $L^2_tH^h_\theta$; the weights below
are applied strip by strip.  All fixed additional angular orders needed
for smooth prescribed boundary data are chosen at the outset.

At this same angular order and in the prescribed relative-height units,
this is the trace-augmented strip convention of \eqref{gauge:norm},
using the real path parameter on each modified band.  The first traces are
controlled on a fixed enlarged strip by \eqref{lin:first-traces},
in the standing regime of a bounded positive diffusion scale and
uniformly comparable values of $E$.  The additional angular derivative
in \eqref{jump:boundary-lift} concerns prescribed outgoing data.

We need a finite-interval inverse more precise than a local strip
estimate.  All the height paths below have length at least a fixed
$T_{\min}>0$: the first fixed-radius exit is separated from the
incoming section, and subsequent exits are farther inward.  We use
overlapping strips $I_\nu$ of lengths between fixed positive
constants, indexed in their order toward the exit.  For any positive
weight $\omega$ with bounded logarithmic rate, put
\[
 \|v\|_{X_\omega}=\sup_\nu\omega_\nu^{-1}X_E^h(v;I_\nu),
 \qquad F_\nu=\|f\|_{L^2(I_\nu;H^h)},\qquad
 \omega_\nu=\inf_{I_\nu}\omega.
\]
For the mean estimate we further require $\omega$ to be
nondecreasing on the long incoming pieces and to vary by a fixed
bounded factor on the modified bands.  Fix its logarithmic-rate bound
and this factor before taking $E_*$ small; the constants
may depend on them.  The weights used below satisfy these
requirements.
With $E_*=\sup E\leq C Y_0^{-2}$, define
\begin{equation}\label{jump:mean-source-norm}
 \mathfrak M_\omega(f)=
 \sup_\nu\frac1{\omega_\nu}
 \left\{\sum_{\mu\leq\nu+1}F_\mu+
 \sum_{\mu>\nu+1}e^{-c(\mu-\nu-1)/E_*}F_\mu\right\}.
\end{equation}
Changing the fixed overlap convention only changes the constants.

\begin{lemma}[Finite height inverse]\label{jump:height-inverse}
On the real parameter interval $[0,T]$, let
$\mathcal L_0$ be the circular operator
\eqref{jump:height-linear} with
\eqref{jump:height-coefficients}, the normalized derivative bounds
just stated, $T\geq T_{\min}$, and $E\leq E_*$ sufficiently small.
For entrance data $\varphi$ in $\mathcal T^D_{E(0),h}$ and outgoing
coordinate-derivative data $\psi$ in $\mathcal T^N_{E(T),h}$,
the problem
\[
 \mathcal L_0v=f,\qquad v(0)=\varphi,\qquad v_s(T)=\psi
\]
has a unique solution in the finite strip domain.  Write subscripts
$m,o$ for its angular mean and nonconstant part.  The mean satisfies
\begin{equation}\label{jump:mean-inverse}
 \|v_m\|_{X_\omega}\leq C\left\{
 \frac{|\varphi_m|}{\omega(0)}
 +\mathfrak M_\omega(f_m)
 +\frac{\|\psi_m\|_{\mathcal T^N_{E(T),h}}}{\omega(T)}
 \right\}.
\end{equation}
There is a weight $\rho_o\asymp (|r|/r_0)^{9/10}$, with
$\rho_o(0)=1$, logarithmic rate $-9/10$ on the long incoming
pieces away from the reserved collars, bounded smooth interpolation
on those $J=1$ collars, and rate zero on the bounded modified bands, for which
\begin{equation}\label{jump:oscillatory-inverse}
 \|v_o\|_{X_{\rho_o}}\leq C\left\{
 \|\varphi_o\|_{\mathcal T^D_{E(0),h}}
 +\sup_\nu(\inf_{I_\nu}\rho_o)^{-1}\|f_o\|_{L^2(I_\nu;H^h)}
 +\frac{\|\psi_o\|_{\mathcal T^N_{E(T),h}}}{\rho_o(T)}
 \right\}.
\end{equation}
The constants are independent of $B,T$ and the exit size
$Y\in[Y_0,cB]$.  A fixed positive separation between the first exit
and the entrance is included in this assertion.
\end{lemma}

\begin{proof}
For the mean put $A=a^{-1}$ and
\[
 H(s,t)=\exp\left(-\int_s^t A(u)\,du\right),\qquad
 F=f_m-q_0v_m .
\]
Solving the equation for $v_{m,s}$ backward from its actual exit gives
the exact formula
\begin{equation}\label{jump:mean-kernel}
 v_{m,s}(s)=H(s,T)\psi_m+
       \int_s^T H(s,t)A(t)F(t)\,dt,\qquad
 \int_s^T H(s,t)A(t)\,dt=1-H(s,T).
\end{equation}
In particular the finite kernel does not have mass one at the
terminal section.  Its missing terminal layer has $L^2$ size
$O(\sqrt{E(T)})$ when it multiplies bounded data.  We keep that layer
in \eqref{jump:mean-kernel}.

Accretivity gives
$|H(s,t)|\leq\exp(-c\int_s^t E^{-1})$ and
$|A(t)|\leq C/E(t)$.  The kernel $H(s,t)A(t)$ has bounded row and
column masses.  For the column estimate one uses the increasing $E$
on the incoming pieces; its variation on the modified bands is
bounded.  Equivalently, bounded logarithmic derivatives and small
$E_*$ give the same estimate on consecutive strips.  Integrating
\eqref{jump:mean-kernel} once yields
\[
 \begin{split}
 v_m(s)={}&\varphi_m+\psi_m\int_0^sH(u,T)\,du
       +\int_0^s m(t)F(t)\,dt+\int_s^T K_+(s,t)F(t)\,dt,\\
 m(t)={}&A(t)\int_0^tH(u,t)\,du,\qquad
 K_+(s,t)=A(t)\int_0^sH(u,t)\,du .
 \end{split}
\]
Here $|m(t)|\leq C$ and
$|K_+(s,t)|\leq C E(s)E(t)^{-1}
\exp(-c\int_s^t E^{-1})$; the latter kernel has row mass
$O(E(s))$.  The past integral has the ordinary Volterra
Gr\"onwall bound $\exp(C\int|q_0|)$, which is uniform, and the
future $q_0v_m$ term is absorbed by taking $E_*$ small.
If $L$ bounds the logarithmic rate of the mean weight, then
\[
 \frac{\omega(t)}{\omega(s)}
 \leq C_{\rm mod}e^{L(t-s)}
 \leq C_{\rm mod}\exp\left(LE_*\int_s^t E^{-1}\right).
\]
Thus $LE_*<c/2$ preserves a fixed part of the fast exponent
in the weighted future and exit kernels.
Schur's inequality for the differentiated formula, followed by
the equation for $a v_{m,ss}$ and the first-trace estimate of
\cref{lin:strip}, proves \eqref{jump:mean-inverse}.  The two
source sums in \eqref{jump:mean-source-norm} are respectively its
cumulative past part and its fast future tail.  The terminal lift
has derivative bulk size $O(\sqrt{E(T)}|\psi_m|)$ and value size
$O(E(T)|\psi_m|)$, so its contribution is precisely the stated
natural Neumann term, uniformly after weighting.

For the nonconstant modes choose $d=(\log\rho_o)_s=-9/10$ on
the long incoming pieces away from the reserved collars and $d=0$ on
the bounded modified bands.  Ramp $d$ smoothly between these values
only on the fixed $J=1$ collars, where the unit incoming estimate for
$c$ holds; in particular $d=0$ before entering each merely accretive
band and until after leaving it.  Their bounded
geometry makes this weight comparable to $(|r|/r_0)^{9/10}$.
Writing $v_o=\rho_o w$ gives
\[
 \rho_o^{-1}\mathcal L_0(\rho_o w)
 =-aw_{ss}+(1-2ad)w_s-cw_{\theta\theta}
             +\{d-a(d^2+d_s)+q_0\}w.
\]
For mean-zero functions the first angular eigenvalue is one.
On the long incoming pieces the real part of the mass therefore has
the strict buffer $1-9/10$ after the small $r_0$ and $E_*$ errors.
On every modified band $d=0$, and $\Re c\geq c_*$ supplies a fixed buffer
because $q_0=O(r^2)$.  On each ramp collar, $J=1$ and
$c=1+O(|r|^2(1+|\log r|)+E)$ while $-9/10\le d\le0$;
the bounded $E d_s$ term is small, so the same fixed incoming buffer
persists there.  Integration against $\bar w$ absorbs the
normalized coefficient derivatives and imaginary drift terms into
$E|w_s|^2$, $|w_\theta|^2$ and this positive mass.  The actual
homogeneous exit is $w_s+dw=0$; its radial and drift boundary
contributions add to
\[
 d\Re a+\tfrac12\Re(1-2ad)=\tfrac12.
\]
Thus the negative incoming weight causes no terminal loss.
Localizing the energy by a smooth weight comparable to
$e^{-\epsilon|s-s_\nu|}$, with $\epsilon$ smaller than the unused
mass buffer, sums the source strips geometrically.  The local
estimate in \cref{lin:strip} then supplies all bulk and first-trace
slots.  Natural boundary lifts prove
\eqref{jump:oscillatory-inverse}.  Finally the finite
Dirichlet--Neumann problem has index zero by its complementing
boundary parametrices and an accretive homotopy at fixed $T$;
the estimates prove injectivity and hence solvability.  This index
calculation is on the ordinary compact $H^2$ domain and $L^2$
range; the commuted estimates and density then give the stated
fixed angular order.
\end{proof}

The outgoing datum in this lemma is the derivative of the
relative-height scalar in the prescribed real coordinate.  Thus
$\psi=\partial_s v=J\partial_t v=r_s\partial_r v$ at the exit;
physical $z$ data are also divided by
$Z$.  At order $h$ its natural norm is
\begin{equation}\label{jump:neumann-trace}
 \|\psi\|_{\mathcal T^N_{E,h}}^2
 =E\sum_{j\in\mathbb Z}(1+j^2)^h
                     (1+\sqrt E|j|)|\psi_j|^2.
\end{equation}
This is the trace norm of \eqref{lin:neumann-trace}, controlled by
the bulk $v_s,\sqrt E\,v_{s\theta},Ev_{ss}$ slots on a terminal
strip.  In particular a derivative trace supplied by another
solution already lies in this space.  For prescribed smooth data
one may use the sufficient lift bound
\begin{equation}\label{jump:boundary-lift}
 \|\operatorname{lift}\psi\|_{X_E^h}
 \leq C\|\psi\|_{\mathcal T^N_{E,h}}
 \leq C\sqrt E\,\|\psi\|_{H^{h+1}_\theta}.
\end{equation}
Only the last, convenient inequality asks for the extra angular
derivative.

The circular residual of \eqref{jump:height-profile} has the required
small size without a spatial analyticity assumption.  Set
$\epsilon_Z=Z^{-2}$.  Its radial equation, multiplied by $r^2$, is
\[
 r^2\left\{\frac{g_c''}{1+Z^2(g_c')^2}
                +(r^{-1}-r/2)g_c'+g_c/2\right\}
 =r^2\epsilon_Z\left\{
 \frac{g_c''}{\epsilon_Z+(g_c')^2}
                       -\frac{\phi''}{(\phi')^2}\right\}.
\]
Since $\phi'\asymp r$, $\phi''=O(1)$,
$d_0'=O(r^{-1})$ and $d_0''=O(r^{-2})$,
the braces are $O(B^{-2}|r|^{-4})$ whenever $|Br|\geq Y_0$.
Thus the residual, with each fixed spatial logarithmic derivative, is
\begin{equation}\label{jump:profile-residual}
 O\bigl(B^{-2}|Br|^{-2}\bigr).
\end{equation}
The same bound holds on the modified bands by the bounded chain-rule factors.

\subsection{The refined exterior contraction}

Fix
\begin{equation}\label{jump:weights}
 \alpha=\frac74,\quad \gamma=\frac9{10},\quad
 \varepsilon=\sigma=\frac1{10},\qquad
 W_o=B^{-\alpha}|r|^\gamma,\quad
 W_m=B^{-2}\bigl(B^{-\varepsilon}+|Br|^{-\sigma}\bigr).
\end{equation}
Changing these four exponents slightly has no effect on the argument.
Write $u=g-g_c$, $u_m=(2\pi)^{-1}\int_{\mathbb S^1}u\,d\theta$
and $u_o=u-u_m$, and use the norm
\begin{equation}\label{jump:refined-norm}
 \|u\|_{\mathcal X}=
 \sup_I\frac{X_E^h(u_m;I)}{W_m(I)}+
 \sup_I\frac{X_E^h(u_o;I)}{W_o(I)}.
\end{equation}
For a path ending at $s=T$, the prescribed remainder derivative is
\[
 \psi=\left.\partial_s(g-g_c)\right|_{s=T}.
\]
For a fixed constant $C_{\rm tr}$, the permitted height data are
exactly the elements of
$\mathcal T^N_{E(T),h}$ satisfying
\begin{equation}\label{jump:permitted-data}
 \|\psi_m\|_{\mathcal T^N_{E(T),h}}\leq C_{\rm tr}W_m(T),
 \qquad
 \|\psi_o\|_{\mathcal T^N_{E(T),h}}\leq C_{\rm tr}W_o(T).
\end{equation}
The families used for first variation are smooth families in this
space.  The canonical prescription $\partial_sg=\partial_sg_c$
has $\psi=0$.  There is no additional $H^{h+1}$ requirement on a
trace obtained from a solution.  We choose $C_{\rm tr}$ below from
the canonical zero-exit estimates, before $B$ and before solving
the interpolated problems.

The incoming mean value is zero by \eqref{jump:matching-Z}; its
oscillatory Dirichlet trace has
$\mathcal T^D_{E(0),h}$ size $O(B^{-1.8})$ and hence fits the
$W_o$ weight.  Lemma~\ref{jump:height-inverse} applies with
$\omega=W_m$ and with a fixed multiple of $\rho_o$ equal, up to
uniform factors, to $W_o$.
The residual \eqref{jump:profile-residual}, propagated by the mean
inverse, fits $B^{-2}|Br|^{-\sigma}$.  There is no oscillatory
reference residual.  Indeed its past strip sum is
$O(B^{-2}|Br|^{-2})$, and the fast future sum has the same bound after
$Y_0$ is fixed large.

We give the entire smallness calculation for the nonlinear map.  In a
fixed multiple of the ball \eqref{jump:refined-norm}, the full height
estimate of \cref{gauge:height} bounds its quadratic and higher
remainder in a strip by
\begin{equation}\label{jump:height-products}
 C\left\{\frac{E^{-1/2}}{|r|^2}(W_m+W_o)^2
       +\frac{E^{-1}}{|r|^4}W_o^2(W_m+W_o)\right\}.
\end{equation}
The difference estimate replaces one weight by the corresponding
difference norm.  Oscillatory output always contains an oscillatory
input.  This last assertion is exact: the operator maps circular
functions to circular functions.  In deriving
\eqref{jump:height-products}, the radial coefficient variation retains
its factor $E$, and the additional radial second-derivative term has
two angular tilts.  At least one differentiated factor stays in strip
$L^2$; the only time-derivative trace loss is $E^{-1/2}$.  Thus no
uncontrolled second-derivative trace is being used.

Put $A=E^{-1/2}|r|^{-2}\asymp B/|r|$; the second coefficient in
\eqref{jump:height-products} is comparable to $A^2$.  The coefficient
factors multiplying a mean or oscillatory error obey
\begin{align}
 AW_m&\leq C\left\{\frac{B^{-\varepsilon}}{|Br|}
                         +|Br|^{-1-\sigma}\right\},
 &AW_o&\leq CB^{-3/4}|r|^{-1/10}.
 \label{jump:lipschitz-factors}
\end{align}
Their logarithmic integrals over $Y_0/B\leq |r|\leq r_0$ are bounded
respectively by
\begin{equation}\label{jump:integrated-factors}
 C(B^{-\varepsilon}Y_0^{-1}+Y_0^{-1-\sigma}),
 \qquad CB^{-13/20}Y_0^{-1/10}.
\end{equation}
The same is true, with smaller bounds, for their squares and products.
The pure oscillatory forcing of the mean is
\begin{align}
 AW_o^2&\leq CB^{-5/2}|r|^{4/5},&
 A^2W_o^3&\leq CB^{-13/4}|r|^{7/10}.
 \label{jump:osc-to-mean}
\end{align}
These powers are logarithmically summable.  Relative to the constant
part $B^{-2-\varepsilon}$ of the mean budget they cost respectively
$O(B^{-1/2+\varepsilon})$ and $O(B^{-5/4+\varepsilon})$.
All other mean terms have an error factor multiplied by one of
\eqref{jump:lipschitz-factors}; all other oscillatory terms retain an
oscillatory factor and have the same small coefficients.  In the
oscillatory inverse the buffer below rate $1$ sums the weighted
propagation kernel.  In the mean inverse use
\eqref{jump:integrated-factors}; its increasing component
$|Br|^{-\sigma}$ is preserved by forward summation.  Consequently the
linear solution operator $T$ satisfies, on each fixed ball,
\begin{equation}\label{jump:contraction}
 \|T\mathcal N(u)-T\mathcal N(v)\|_{\mathcal X}
 \leq C\bigl(B^{-\delta}+Y_0^{-\delta}\bigr)
                           \|u-v\|_{\mathcal X}
\end{equation}
for a fixed $\delta>0$.  Constants absorb the bounded bands.
The linear image bounds used next are uniform once $Y_0$ exceeds
a fixed threshold.
First take $\psi=0$ and choose a fixed ball from the reference
residual and incoming lift.  For $Y_0,B$ above fixed preliminary
thresholds, \eqref{jump:contraction} gives a canonical solution
with a bound uniform for all larger thresholds and independent
of $C_{\rm tr}$.  Its natural derivative trace at any later common
section is bounded by a fixed multiple of $W_m,W_o$, by
\eqref{jump:neumann-trace}.  Choose $C_{\rm tr}$ larger than
these two canonical trace constants.  The reference residual,
incoming lift and the entire class \eqref{jump:permitted-data}
are then bounded in the same norm.  Choose a fixed
$C_{\rm ref}$ large enough to contain their images and run the
contraction on
\begin{equation}\label{jump:refined-ball}
 \mathfrak B_{\rm ref}(C_{\rm ref})
       =\{u:\|u\|_{\mathcal X}\leq C_{\rm ref}\}.
\end{equation}
The estimates also hold on twice this fixed ball, with contraction
constant less than $1/2$ after enlarging $Y_0$ for these fixed
ball constants and then taking $B$ large.  The
contraction produces the unique exterior solution in that ball
with its stated entrance and exit data, and
\begin{equation}\label{jump:refined-bounds}
 X_E^h(g_o;I)\leq CW_o,
 \qquad X_E^h(g_m-g_c;I)\leq CW_m.
\end{equation}

The same estimates give an inhomogeneous-entrance difference bound.
If two solutions for the same path and parameter lie in twice
\eqref{jump:refined-ball}, their
mean-value linearization obeys \eqref{jump:contraction}.  Subtracting
the equations and absorbing that term gives
\begin{equation}\label{jump:refined-difference}
 \|u-\widetilde u\|_{\mathcal X}\leq C\left\{
 \frac{\|\varphi_m-\widetilde\varphi_m\|_{\mathcal T^D}}{W_m(0)}
 +\frac{\|\varphi_o-\widetilde\varphi_o\|_{\mathcal T^D}}{W_o(0)}
 +\frac{\|\psi_m-\widetilde\psi_m\|_{\mathcal T^N}}{W_m(T)}
 +\frac{\|\psi_o-\widetilde\psi_o\|_{\mathcal T^N}}{W_o(T)}
 \right\}.
\end{equation}
The trace norms here use their respective entrance or exit diffusion
scales.  This estimate will propagate the small, nonzero entrance
shift caused by global normalization.

For each use of holomorphy, hold the central contour and its
phase-alignment profile fixed while $p$ varies on its parameter ball.
All these maps are holomorphic in the finite parameters when their
reference coefficients and prescribed trace data are varied
holomorphically on the prescribed real intervals, and the trace
families satisfy \eqref{jump:permitted-data} uniformly on the
outer ball.  For this purpose identify their trace spaces with
the fixed central-parameter norm; the varying diffusion scales
are uniformly comparable there.  Use the outer and inner balls
$\mathcal B_s^{\rm out}$ and $\mathcal B_s$ fixed in
\eqref{arr:parameter-reserve}, choosing their constants there
small enough to keep $Z$, the entrance data and the strict
coefficient margins in the same regimes.  Uniform contraction
gives holomorphic dependence on $\mathcal B_s^{\rm out}$.
Every finite-parameter derivative assertion below is on
$\mathcal B_s$, at distance at least
$(c_{\rm out}-c_{\rm in})|s|^L$ from the outer boundary.
Cauchy estimates there
cost only a fixed power for each fixed derivative order.  The first
logarithmic-exit and terminal-interpolation derivatives used below
instead follow from the pulled equations.

\subsection{Exponential shielding and normalized gluing}

The exterior solve has so far had prescribed entrance and exit data.
We now turn it into a normalized full-surface solution, keeping the
numbers $Z$ and $Z_*$ from the preliminary match fixed.  Write
$(V_{\rm pre},\lambda_{\rm pre})$ for the common normalized
preliminary graph tuple and its compact multiplier, with the common
nonadverse problems fixed above.  For each height
path and permitted $\psi$, call the exterior solution with entrance
exactly equal to that of $V_{\rm pre}$ the canonical exterior
solution.  Scalar graph tuples are always compared in the same
prescribed real charts.

On the common real interval $[0,T_{\rm pre}]$, before the phase
alignment or either turn, subtract the
preliminary solution from the canonical exterior solution.  Their difference
has zero entrance value and solves a homogeneous linear difference
equation; the coefficients satisfy the same height smallness
conditions.  Indeed, there $E\asymp B^{-2}$ and the relevant
strip sizes are at most $C(B^{-7/4}+B^{-1.8})$, so the shifted
coefficient smallness required by \cref{lin:shielding} follows from
$B(B^{-7/4}+B^{-1.8})=o(1)$.  Choose
$d=c_1B^2$ with $Ed$ a sufficiently small fixed constant, and write
$v=e^{d(t-T_{\rm pre})}w$.  For the model principal part
$\partial_t-E\partial_t^2-A_2\partial_\theta^2$, the conjugated
operator is
\begin{equation}\label{jump:large-shift}
 -E\partial_t^2+(1-2Ed)\partial_t
 -A_2\partial_\theta^2+d-Ed^2.
\end{equation}
The positive gain is $d-O(Ed^2)$, comparable to $d$.
The true transformed outgoing condition is
$w_t+dw=f$; its fast-root denominator is
$(\lambda_f-d)+d=\lambda_f$.  Thus no small denominator or
exponential loss is introduced.  For real frozen coefficients its
terminal energy coefficient is exactly
$Ed+(1-2Ed)/2=1/2$.  The accretive perturbation and trace estimates
in \cref{lin:shielding} preserve this conclusion.  Alternatively,
cut off toward the far side of the overlap and use its Dirichlet
trace.  In either implementation, every needed fixed derivative on
a smaller inner joining band satisfies
\begin{equation}\label{jump:shield}
 \|v\|_{\rm join}\leq CB^M e^{-c_{\rm sh}B^2}.
\end{equation}
Here the joining norm includes the fixed spatial derivatives needed
for the splice.  The joining bands are fixed, so $c_{\rm sh}>0$ is
uniform in all exit sizes $Y\in[Y_0,cB]$ and in the permitted data.

Splice on that band, obtaining a graph tuple $V_{\rm app}$ which
equals $V_{\rm pre}$ on the compact observation region and the
canonical exterior outside the joining band.  Use
$\lambda_{\rm pre}$ as its multiplier.  The resulting augmented
residual has the size in \eqref{jump:shield}, and the compact
observations remain unchanged.

We verify the stopped index hypothesis of \cref{lin:global} for
this spliced linearization.  First use smooth permitted data, as
in the action families.  On the refined exterior,
\eqref{jump:height-coefficients} and the projected height
smallness give strict row-normalized accretivity.  The
exponentially small splice preserves it on the fixed joining
bands, and the compact and conical pieces are the preliminary
ones.  The prescribed real charts retain nonzero graph-speed
factors, and the correction exit is the pure relative-height
derivative.  At every finite stop we may therefore use the
same graph-speed conjugations, patched nonzero row multiplier,
exit-collar Robin removal, and convex principal-tensor homotopy
as in the proof of \cref{arr:preliminary-solves}.  The symbols
remain elliptic and the exits complementing.  Retaining the
compact multiplier columns and the equal-dimensional
observation augmentation gives index zero on the fixed
compact domain.  This verifies the index hypothesis for the
spliced operator; no estimate along that artificial homotopy
is used.

The global inverse of \cref{lin:global} is applied with cylindrical
growth rate $2.2$, strictly above the propagation rate $2.1$, bounded transitions on the fixed
scaled bands, and a large fixed height growth rate.  On the height
interval of length $O(\log B)$ its weight costs only a power of $B$.
Keeping those weights until after the global estimate is essential:
the forward estimate with a strict buffer prevents a normalized
approximate kernel from escaping to the far end, so its nonzero
compact limit would contradict the Gaussian normalized inverse.
The cited global theorem supplies this estimate and surjectivity by
the index-zero argument on compact stopped problems with complementing
exits, followed by uniform conical exhaustion.  In physical graph norms
its inverse therefore costs at most $CB^q$.
For a nonsmooth datum in the full class
\eqref{jump:permitted-data}, approximate it by smooth data in
the natural trace norm within the same fixed ball.  The
refined difference estimate gives convergence of the exterior
graphs, and \cref{gauge:height} gives convergence of their
linearizations in the domain-to-range operator norm at each
fixed stop.  A small Neumann perturbation of the uniform
smooth-data inverses, followed by density, supplies the same
inverse for that datum.  The actual action families already
have smooth comparison data or smooth interior traces.

This inverse now corrects an exponentially small error.  Reserve
once and for all a strict exponent margin
$0<c_{\rm gl}<c_{\rm sh}$.  In the affine domain with the prescribed
outgoing derivatives, define the augmented uniqueness ball
\begin{equation}\label{jump:augmented-ball}
 \mathfrak U(V_{\rm app})=
 \left\{(V,\lambda):
 \|(V,\lambda)-(V_{\rm app},\lambda_{\rm pre})\|_{X_B}
                          <e^{-c_{\rm gl}B^2}\right\}.
\end{equation}
The norm $X_B$ is that of \cref{lin:global}; in particular it
includes the Euclidean multiplier norm.  The polynomial losses of
the inverse and of the nonlinear map give a contraction constant
at most $CB^m e^{-c_{\rm gl}B^2}$ on this ball, for a fixed $m$.
The image of its center is $CB^M e^{-c_{\rm sh}B^2}$, which lies
strictly inside the ball for large $B$.  The augmented contraction
therefore has a unique normalized solution there, with correction
$(k,\ell)=(V-V_{\rm app},\lambda-\lambda_{\rm pre})$ satisfying
\begin{equation}\label{jump:correction-bound}
 \|(k,\ell)\|_{X_B}\leq CB^M e^{-c_{\rm sh}B^2}.
\end{equation}
Increasing $M$ gives the same estimate in the unweighted graph
and first-trace norms, including conversion to physical height.
This is a second contraction, distinct from
\eqref{jump:contraction}; the polynomial global inverse has only
been applied after shielding.  The corrected solution satisfies
\eqref{jump:refined-bounds} up to errors smaller than every power
of $B$.  For data families holomorphic in $p$ in the fixed identified
trace spaces, the canonical exterior and splice center are holomorphic,
and the analytic augmented contraction gives a holomorphic normalized
solution on the outer parameter ball.  Cauchy differentiation for
these families on the fixed inner parameter ball
only increases $M$ in \eqref{jump:shield} and
\eqref{jump:correction-bound}.

The following comparison puts the solutions used in the action
calculation into the two uniqueness classes just constructed.

\begin{proposition}[Compatibility of the normalized height families]
\label{jump:compatibility}
Fix the common parameter $p$ on the inner parameter ball and the
matched numbers $Z,Z_*$ on each adverse end.  Choose the fixed
constants in the preliminary and refined contraction balls before
$B$.  Then the following hold for large $B$.
\begin{enumerate}
 \item At a fixed positive physical radius on the first height
 turn, the normalized height solution with derivative interpolated
 between the preliminary circular profile and $g_c$ belongs to the preliminary uniqueness
 ball of \cref{arr:preliminary-solves} and to a fixed refined
 ball.  It is the preliminary normalized solution with those
 same path and exit data.
 \item Let $V_L$ be either of the two normalized solutions on
 longer canonical paths with zero remainder derivative at its own exit,
 and let $g_L$ be its relative height.
 At a common shorter incoming height section let
 $\psi_L=\partial_s(g_L-g_c)$ be its actual remainder derivative,
 in the shorter scalar coordinate.  The restriction
 $(V_L|_{\rm short},\lambda_L)$ belongs to the augmented ball
 \eqref{jump:augmented-ball} for the shorter canonical exterior
 solve with datum $\psi_L$.  It equals that shorter normalized
 solve.  The linear interpolation between the two data
 $\psi_-$ and $\psi_+$ gives an actual normalized family whose
 endpoints are exactly the two restrictions.
\end{enumerate}
All constants are uniform in the shorter size
$Y\in[Y_0,cB]$.  The comparison retains the compact multiplier,
and all members use the same $Z,Z_*$ and observation value.  The
nonadverse paths, exit sections, and outgoing derivative prescriptions
are the common frozen ones specified before the height construction.
\end{proposition}

\begin{proof}
For the first assertion take the first turn to end at
$|r|=r_1>0$, with $r_1$ fixed below $r_0$.  Its logarithmic
length is bounded and $E\asymp B^{-2}$.  On its terminal band the
preliminary comparison is circular.  In relative-height units
it is
\[
 g_R=(Z_0/Z)\phi,\qquad
 g_c-g_R=(1-Z_0/Z)\phi+Z^{-2}d_0.
\]
The last difference and every fixed logarithmic derivative are
$O(B^{-1.8})$ on this fixed-radius interval.  In particular the
known smooth comparison derivative changes by $O(B^{-1.8})$ in
$H^{h+1}$, and its natural lift has size $O(B^{-2.8})$ by
\eqref{jump:boundary-lift}.  This assertion concerns the smooth
comparison profiles, not an unknown solution trace.
After the fixed graph-unit conversion, division by the preliminary
weight $\rho_B(T)\asymp B^{2.2}$ gives $O(B^{-5})$, inside
\eqref{arr:allowed-traces} throughout this interpolation.
For each fixed path, $Z$ is the bounded mean trace of the preliminary
holomorphic solution at the fixed incoming section.  The chosen branch
of $Z_0$ and the formulas for $g_R$ and $g_c$ are therefore holomorphic
in $p$ on the outer ball.  Their comparison derivatives, evaluated at
the prescribed real path parameter at the exit, are holomorphic in the
identified trace space, as is their affine interpolation.  This verifies the
parameter-data condition in \cref{arr:preliminary-solves}.

Apply \cref{jump:height-inverse} on this bounded interval with the
constant total budget $W_{\rm fix}=B^{-1.8}$.  The incoming
oscillation is $O(W_{\rm fix})$, its mean is zero, the reference
residual is $O(B^{-4})$, and the preceding boundary lift is
$O(B^{-2.8})$.  Here $E^{-1/2}|r|^{-2}\asymp B$, so the full
nonlinear remainder and Lipschitz factor are
\[
 O(BW_{\rm fix}^2+B^2W_{\rm fix}^3),\qquad
 O(BW_{\rm fix}+B^2W_{\rm fix}^2)=O(B^{-0.8}).
\]
There is consequently a unique fixed-interval solution in
$\sup_I X_E^h(u;I)\leq C_{\rm fix}B^{-1.8}$ for a fixed
$C_{\rm fix}$.  Project its equation onto the mean and use the
zero incoming mean.  The same inverse and the projected products
give the stronger bound
\[
 \sup_I X_E^h(u_m;I)
 \leq C\{B^{-4}+B^{-2.8}
              +BW_{\rm fix}^2+B^2W_{\rm fix}^3\}
 =O(B^{-2.6}).
\]
At fixed $r$, $W_o\asymp B^{-7/4}$ and
$W_m\asymp B^{-2.1}$.  Thus its oscillation and mean lie in the
refined ball for large $B$.  It has the same entrance and exit
data as the canonical refined exterior, so uniqueness in
\eqref{jump:refined-ball} identifies those two exterior solves.
The normalized correction
\eqref{jump:correction-bound} preserves this membership.

For clarity, the preliminary ball here is the growth-weighted
$X_B$ ball of radius $C_{\rm pre}B^{-4}$ about the comparison pair
of \cref{arr:preliminary-solves}, after its prescribed boundary
lift.  The weight on this fixed outer band is comparable to
$B^{2.2}$, so $B^{-1.8}/B^{2.2}=B^{-4}$.  The core and conical
parts equal those of $V_{\rm pre}$ before the new normalized
correction.  There $V_{\rm pre}$ differs from the shared preliminary
comparison pair by $O(B^{-4})$ by \cref{arr:preliminary-solves};
\eqref{jump:correction-bound} controls the additional discrepancy
and the multiplier.  Choose $C_{\rm pre}$ larger than these fixed
bounds before taking $B$ large.  The preliminary contraction
remains a contraction on this fixed enlarged ball because its
outer Lipschitz factor is $O(C_{\rm pre}B^{-0.8})$.  The new
normalized solution and the preliminary solution solve the same
augmented equation, observations and exit problem inside that
ball.  Its uniqueness identifies them.  This proves the first
assertion and connects the fixed-radius end of the array
homotopy to the refined family.

For the second assertion, the natural trace theorem applied on an
enlarged strip of each longer canonical solution gives
\[
 \|\psi_{L,m}\|_{\mathcal T^N}\leq C W_m(T),\qquad
 \|\psi_{L,o}\|_{\mathcal T^N}\leq C W_o(T).
\]
Its normalized correction adds only
$B^M e^{-c_{\rm sh}B^2}$ in these norms.  The earlier choice
of $C_{\rm tr}$ includes these two canonical trace families
for large $B$.  Thus
\eqref{jump:permitted-data} holds for their actual traces and
for their linear interpolations.
For fixed paths and shorter sections, the longer zero-remainder-exit
solutions are holomorphic in $p$ by the preceding contractions.
After the fixed real pullbacks, restriction and the natural trace map
are bounded complex-linear maps into the identified shorter trace
space.  Thus $\psi_-(p)$, $\psi_+(p)$, and their affine interpolation
are holomorphic there.  Variation of the contour or shorter section
uses the separately proved smooth pulled-domain derivatives.

Let $g_{\rm short}^{\rm can}$ be the shorter canonical exterior
with entrance $g_{\rm pre}$ and datum $\psi_L$.  The restriction
$g_L$ satisfies the same exterior equation because the compact
multiplier sources vanish on this height region.  It has the same
actual exit derivative, but its entrance differs
by
\[
 e_0=(g_L-g_{\rm pre})|_{s=0},\qquad
 \|e_0\|_{\mathcal T^D}\leq CB^M e^{-c_{\rm sh}B^2},
\]
by \eqref{jump:correction-bound} and the trace estimate.  Both
exteriors lie in twice the refined ball.  Applying
\eqref{jump:refined-difference} to their mean-value difference,
with zero exit discrepancy and this nonzero entrance discrepancy,
gives
\begin{equation}\label{jump:restriction-difference}
 \|g_L-g_{\rm short}^{\rm can}\|_{\mathcal X}
             \leq CB^{M'}e^{-c_{\rm sh}B^2}.
\end{equation}
Dividing by $W_m(0)$ and $W_o(0)$ has only increased the fixed
power.  No zero-entrance shielding has been applied to this
difference.

On the core the longer restriction differs from $V_{\rm pre}$ by
its correction \eqref{jump:correction-bound}, including
$\lambda_L-\lambda_{\rm pre}$.  On the shorter exterior it differs
from $g_{\rm short}^{\rm can}$ by
\eqref{jump:restriction-difference}.  On the joining band the
remaining comparison between $g_{\rm short}^{\rm can}$ and
$g_{\rm pre}$ has zero entrance and is exactly the canonical
comparison to which \eqref{jump:shield} applies.  Combining these
three estimates in the prescribed charts and in $X_B$ costs only
another fixed power, and hence
\[
 \|(V_L|_{\rm short},\lambda_L)
       -(V_{\rm app,short},\lambda_{\rm pre})\|_{X_B}
 \leq CB^{M''}e^{-c_{\rm sh}B^2}
                  <e^{-c_{\rm gl}B^2}.
\]
The restriction is therefore in the shorter augmented ball.
It has the same PDE, observation values and outgoing derivative
as the normalized shorter solution; uniqueness in
\eqref{jump:augmented-ball} identifies them.

Finally put $\psi_\xi=(1-\xi)\psi_-+\xi\psi_+$ for
$0\leq\xi\leq1$, simultaneously on all adverse ends.
The natural trace balls in \eqref{jump:permitted-data} are convex.
The exterior contraction and augmented correction therefore give
a smooth normalized family with these data.  Differentiating the
exterior contraction in $\xi$ and using the lift of
$\psi_+-\psi_-$ gives a uniform bound for
$\partial_\xi u$ in the refined norm; its entrance derivative
is zero.  Shielding this canonical derivative on the fixed
joining band and then differentiating the augmented correction
gives \eqref{jump:correction-bound} with a larger fixed power.
At a fixed section $y=Y$, conversion to the scaled height gives
\begin{equation}\label{jump:interpolation-height}
 \|\partial_\xi h(Y,\cdot)\|_{H^h}
 \leq C\{B^{-\varepsilon}+Y^{-\sigma}
                   +B^{-13/20}Y^{9/10}\}
       +B^M e^{-c_{\rm sh}B^2}.
\end{equation}
Indeed $h=(Z_*-Zg)Z_*/2$ and $|ZZ_*|\asymp B^2$, so this is
exactly the refined derivative bound multiplied by $B^2$.
The endpoint identification just proved makes this the required
family between the two restrictions.
\end{proof}

\subsection{Endpoint errors in the action scale}

The estimates have two consequences for which polynomial accuracy
in physical coordinates alone would be insufficient.

\begin{lemma}[Variation of a logarithmic exit]\label{jump:moving-exit}
Keep $p,Z,Z_*$ fixed and use the canonical prescription
$\partial_s(g-g_c)=0$ at the moving exit.  Put $\eta=\log Y$.
For the normalized family obtained above, its physical boundary
position $X_{\rm e}$ satisfies
\begin{equation}\label{jump:exit-derivative}
 \|\partial_\eta X_{\rm e}\|_{H^h_\theta}
                 \leq B^{-1}P(Y),\qquad Y_0\leq Y\leq cB,
\end{equation}
where $P$ is a fixed polynomial independent of $B,Y$.  The assertion is uniform on the inner
parameter ball.
\end{lemma}

\begin{proof}
Near any chosen $\eta$, pull neighboring paths to its fixed real
parameter interval by maps $\Psi_\eta$ which are the identity
before the last bounded bands.  Their derivatives in the path
coordinate and in $\eta$ are bounded in logarithmic units.
This is possible for the homothetic bypasses and their preceding
shortening bands.  First let $g_\eta$ denote the canonical exterior
before normalization, and put
$u_\eta=(g_\eta-g_c)\circ\Psi_\eta$ and use dots for derivatives
at a fixed pulled coordinate.  The exact remainder problem has
the form
\[
 \mathcal L_\eta u_\eta+\mathcal N_\eta(u_\eta)=-R_\eta,\qquad
 u_\eta(0)=u_{\rm in},\qquad (u_\eta)_s(T)=0.
\]
The incoming value is independent of $\eta$.  The pulled exit
condition is still homogeneous because it is a coordinate
derivative of the remainder.  The locally parameterized normalized
branches are smooth by the uniform contractions for these smooth
coefficient maps.  They agree on overlaps by uniqueness in
\eqref{jump:augmented-ball}, so they describe one smooth shortening
family.  Differentiation therefore gives
\begin{equation}\label{jump:differentiated-exit-equation}
 \begin{split}
 \{\mathcal L_\eta+D\mathcal N_\eta(u_\eta)\}\dot u
   &=-\dot R_\eta-(\dot{\mathcal L}_\eta)u_\eta
                    -(\partial_\eta\mathcal N_\eta)(u_\eta),\\
 \dot u(0)&=0,\qquad \dot u_s(T)=0 .
 \end{split}
\end{equation}
No second-derivative boundary trace appears.

The explicit coefficient variations are supported on the last
bounded bands.  Differentiating the full pulled drift and then
dividing it as above gives, schematically, the source slots
\[
 O(1)u_s,\quad O(E)u_{ss},\quad O(1)u_{\theta\theta},
 \quad O(\sqrt E)u_{s\theta},\quad O(1)(u,u_\theta).
\]
Their coefficients are bounded by a fixed polynomial in $Y$,
with no independent power of $B$.  These are precisely the normalized
bulk slots; a background Hessian stays in its bulk factor.
The circular coefficient variations preserve the mean and
nonconstant projections.  Moreover
$\dot R_\eta=O(B^{-2}|Br|^{-2})$ up to such a polynomial,
and the differentiated nonlinear coefficients obey the
projected products in \eqref{jump:height-products}.

Apply the two parts of \cref{jump:height-inverse} to
\eqref{jump:differentiated-exit-equation} and absorb
$D\mathcal N_\eta$ with \eqref{jump:contraction}.
The explicit mean sources occupy only a bounded number of
terminal strips, so their cumulative term in
\eqref{jump:mean-source-norm} contributes no factor $\log B$.
The oscillatory estimate uses its strict angular buffer.
It follows that
\begin{equation}\label{jump:log-remainder-bound}
 X_E^h(\dot u_m;I)\leq P(Y)W_m(I),\qquad
 X_E^h(\dot u_o;I)\leq P(Y)W_o(I).
\end{equation}

At the exit $|r|\asymp Y/B$ and $|\dot r|\leq C|r|$.  Since
$Z$ is fixed,
\[
 |Z\,\partial_\eta(g_c\circ\Psi_\eta)|
 \leq C\{B|r|^2+B^{-1}(1+|r|^2|\log r|)\}
 \leq B^{-1}P(Y).
\]
The first inequality uses
$r\partial_r\phi=O(r^2)$ and
$r\partial_rd_0=O(1+r^2|\log r|)$.
The first traces in \eqref{jump:log-remainder-bound} give
\[
 |Z\dot u|\leq CBP(Y)(W_m+W_o)
 \leq B^{-1}P(Y),
\]
because, up to fixed comparison constants,
$BW_m\leq CB^{-1}(B^{-\varepsilon}+Y^{-\sigma})$
and
$BW_o\leq CB^{1-\alpha-\gamma}Y^\gamma
       =CB^{-33/20}Y^{9/10}$.
A derivative trace used for a slope can cost an additional
factor $E^{-1/2}\asymp Y$, which is also included in $P$.
Together with $|\dot r|\leq CB^{-1}Y$, these are the desired
physical scales for the canonical exterior.

It remains to check the global correction.  On the fixed overlap
before the support of the pulled coefficient variations, the
canonical derivative
$\partial_\eta(g_\eta-g_{\rm pre})$ has zero entrance value and
satisfies the homogeneous linearized equation.  Its outer size
from \eqref{jump:log-remainder-bound} is at most a fixed power
of $B$, since $Y\leq cB$.  Applying \cref{lin:shielding} to
this zero-entrance derivative gives
$B^M e^{-c_{\rm sh}B^2}$ on the joining band.  Thus the derivative
of the spliced residual has that size.  This step uses the
canonical local derivative, not the nonzero-entrance
restriction difference in \eqref{jump:restriction-difference}.

Differentiate the augmented correction equation at fixed
observations and with homogeneous pulled derivative data for
the correction.  The unknown includes $\dot\ell$, the multiplier
derivative.  Its inverse and the coefficient derivatives cost
only fixed powers in the global norms, while the differentiated
residual and every term containing the correction have the
exponent in \eqref{jump:correction-bound}.  Hence, after
increasing $M$,
\[
 \|(\dot k,\dot\ell)\|_{X_B}
       +\|\partial_\eta X_{k,{\rm e}}\|_{H^h}
       \leq CB^M e^{-c_{\rm sh}B^2}\leq B^{-1}.
\]
The last inequality is uniform for $Y\in[Y_0,cB]$.
Adding this correction proves \eqref{jump:exit-derivative}.
\end{proof}

\begin{lemma}[Moving exits and the exterior limit]\label{jump:endpoints}
For the common parameter $p$, first make the high-stop synchronization
\eqref{jump:nonadverse-sync}, then shorten all adverse turns
from fixed scaled size to $|y|\asymp Y_0$.  At each fixed sufficiently
large $Y_0$ and for large $B$, this changes either original action by at most
\begin{equation}\label{jump:endpoint-bound}
 C e^{-cY_0^2}\sum_{j\ \mathrm{adverse}}
 B^{-2}\exp\bigl(-\Re Z_{*,j}^2/4\bigr).
\end{equation}
At each fixed sufficiently large $Y_0$, as $B\to\infty$ the height
functions in \eqref{jump:scaling} have smooth subsequential limits
on compact parts of the two exterior paths and of the interpolated
incoming problems.  Their oscillations vanish, and their radial
limits obey
\begin{equation}\label{jump:translator}
 h'=l,\qquad l'=(1+l^2)(1-l/y).
\end{equation}
On a homothetic bypass of size $Y$ these limits satisfy
\begin{equation}\label{jump:translator-matching}
 h=y^2/2-\log y+O(Y^{-\sigma}),\qquad
 l/y=1+O(Y^{-\sigma}),
\end{equation}
with harmlessly reduced $\sigma>0$, uniformly also in the
interpolation parameter.
\end{lemma}

\begin{proof}
At $r=2y/Z_*$, the oscillatory height error is bounded by
\[
 CB^2W_o\leq CB^{2-\alpha-\gamma}|y|^\gamma
             =CB^{-13/20}|y|^{9/10},
\]
and the mean error by $C(B^{-\varepsilon}+|y|^{-\sigma})$.
Insert \eqref{jump:dstar} into
\eqref{jump:height-profile}.  The constant terms cancel because
of the definition of $Z_*$, giving
\begin{equation}\label{jump:h-match}
 h-(y^2/2-\log y)
   =O(|y|^{-\sigma})+o_B(1)
\end{equation}
for fixed $y$.  The derivative traces may cost a further factor
$|y|$, which is harmless for the relative slope in
\eqref{jump:translator-matching}.  Fixed higher spatial regularity follows
by differentiating the equation and the strip estimates; on each
fixed exterior annulus the scaled coefficients are uniformly
elliptic.  The interpolation estimates in
\cref{jump:compatibility} give the same conclusion for the
interpolated problems.

One can also identify the limiting equation before taking a limit.
With $D_yh=(h_y,h_\theta/y)$, substitution of
\eqref{jump:scaling} into \eqref{jump:height-equation} gives
\begin{equation}\label{jump:exact-scaled-equation}
 \left(I-\frac{D_yh\otimes D_yh}{1+D_yh\cdot D_yh}\right)
       :D_y^2h
       =1+\frac{2(yh_y-h)}{Z_*^2},
\end{equation}
where $D_y^2h$ is the polar Euclidean Hessian.  Oscillations vanish
by the preceding bound.  Letting $B\to\infty$ in
\eqref{jump:exact-scaled-equation} gives
$h''/(1+(h')^2)+h'/y=1$, namely \eqref{jump:translator}.
This uses only regions $|y|\geq Y_0$.

For the endpoint estimate retain the full $\phi$ in
\eqref{jump:height-profile}.  At a complex exit of size
$Y\asymp|Br|$, its negative $\Re(Z^2r^2)$ and the identity
$\phi^2=1-r^2/2$ give
\begin{equation}\label{jump:gaussian-gap}
 \Re(z^2-Z_*^2)\geq cY^2-C(1+\log Y).
\end{equation}
Here the logarithmic correction is in $\log(Br)$, not $\log B$:
\eqref{jump:scaling} subtracts exactly the latter constant.
The residual errors in \eqref{jump:refined-bounds} are bounded in
scaled height by $C(1+Y^\gamma)$, uniformly in $B$, and so can be
absorbed in the right side of \eqref{jump:gaussian-gap} for large
$Y_0$.  The quadratic errors from replacing $Z$ by $Z_*$ are
smaller still.  On the initial fixed-radius pieces any polynomial
in $B$ is absorbed by the extra $e^{-cB^2}$.

By the first part of \cref{jump:compatibility}, the canonical
normalized family at fixed radius is the same array family after
the preliminary deformation and comparison-data interpolation.
That initial deformation costs a fixed power of $B$ times
$e^{-\Re Z_0^2/4-cB^2}$.  The relations
$Z/Z_0=1+O(B^{-1.8})$ and \eqref{jump:scaling} give
$Z_*^2-Z_0^2=o(B^2)$, so this cost is negligible in the
scale $B^{-2}e^{-\Re Z_*^2/4}$.
The earlier synchronization contributes
\eqref{jump:nonadverse-sync}.  The fixed choice of $D$ makes this
$o(B^{-2}e^{-\Re Z_{*,j}^2/4})$ for every adverse $j$; at each fixed
$Y_0$ it is absorbed in \eqref{jump:endpoint-bound} for large $B$.

Along the subsequent shortening, \eqref{jump:exit-derivative}
bounds the physical boundary displacement by $B^{-1}P(Y)$.
The boundary circumference is $O(B^{-1}Y)$, and slopes,
conormals and their continued area branches have at most
polynomial size in $Y$.  Indeed the first traces give
$z_r=O(Y+1+Y^\gamma)$ and
$z_\theta/r=O(B^{-13/20}Y^{-1/10})$ at the exit, while the
relative slope smallness in \cref{gauge:height} keeps the
continued denominators in the same branches as the circular
profile.
Thus both the displacement and momentum terms in first variation
are bounded by
\[
 CB^{-2}P(Y)\exp(-\Re Z_*^2/4-cY^2)
\]
per unit logarithmic size, for a fixed polynomial $P$.  The
observation terms vanish because $p$ is fixed.  Integrating
over $Y_0\leq Y\leq CB$ and decreasing $c$ absorbs $P$ and proves
\eqref{jump:endpoint-bound}.  The nonadverse endpoints stay at
$\Re z^2=DB^2$.  Their derivative prescriptions and paths are fixed,
but their boundary heights can vary through the normalized correction;
their fixed-exit momentum terms remain in the first variation and are
bounded by $C_D B^M e^{-DB^2/5}$ using its polynomial derivative bounds.
The same choice of $D$ absorbs them in the asserted scale.  This proves all
claims, with $B\to\infty$ at fixed $Y_0$ before $Y_0$ is increased.
\end{proof}

\subsection{The scalar Stokes constant}

We supply the ODE argument, including its divergence test, since
the nonvanishing of the eventual action coefficient depends on it.
Equation~\eqref{jump:translator} is the radial translating-soliton
equation in dimension two.  Altschuler--Wu construct the real entire
rotational translator \cite[Corollary~3.3]{AW1994}; its real solutions and quadratic--logarithmic
asymptotics are studied in \cite[Section~2, equation~(2.1) and
Lemmas~2.1--2.2]{CSS2007}.  The complex lateral solutions and their
nonzero difference are constructed here.  General nonlinear Stokes
theory provides related ODE methods \cite{Costin1998}; no summability
or nonvanishing theorem is imported for this equation.

\begin{lemma}[A nonzero lateral difference]\label{jump:stokes}
Equation \eqref{jump:translator} has two exterior lateral slope
solutions selected by $l/y\to1$ on sectors reaching
$\arg(y^2)=\pm2\pi/3$.  On their common positive ray,
\begin{equation}\label{jump:stokes-difference}
 l_+(y)-l_-(y)=c_0y^3e^{-y^2/2}(1+O(y^{-2})),
 \qquad c_0\ne0.
\end{equation}
Finite homothetic bypass solutions satisfying
\eqref{jump:negative-exit} and
\eqref{jump:translator-matching} have the same difference at
$y=Y$, up to $Y^3e^{-Y^2/2}O(e^{-cY^2})$.
\end{lemma}

\begin{proof}
Put $\tau=y^2/2$ and $q=1-l/y$.  The exact equation is
\begin{equation}\label{jump:q-equation}
 q_\tau+q=2q^2-q^3+\frac{1-2q}{2\tau}.
\end{equation}
Choose $\alpha_0>0$ and
$2\pi/3<\vartheta<\pi$ with
$\vartheta+\alpha_0<\pi$, and put $v=e^{i\vartheta}$.
On the translated cone
\[
 \Omega_+(R)=
 \{a e^{-i\alpha_0}+bv:a>R,\ b>0\}
\]
the rays $\tau+tv$, $t\geq0$, remain in the cone and
$|\tau+tv|\geq c(|\tau|+t)$.  As $-\Re v>0$, the map
\begin{equation}\label{jump:volterra}
 (\mathcal Tq)(\tau)=-v\int_0^\infty e^{tv}
 \left(2q(\tau+tv)^2-q(\tau+tv)^3+
                 \frac{1-2q(\tau+tv)}{2(\tau+tv)}\right)dt
\end{equation}
is a contraction on a fixed sufficiently large ball in
$\sup|\tau q(\tau)|$ when $R$ is large.  Indeed its inhomogeneous
part is $O(|\tau|^{-1})$, and its derivative in $q$ is
$O(R^{-1})$ there.  Its holomorphic fixed point satisfies
\eqref{jump:q-equation}, by differentiating under the integral
and integrating once in $t$.  It covers every closed angular
subsector of $(-\alpha_0,\vartheta)$ at sufficiently large
radius.  Reflection gives $q_-$ below.  Uniqueness in these
classes is also immediate: the difference of two small solutions
satisfies $\delta'=(-1+O(\tau^{-1}))\delta$, and a nonzero such
difference grows exponentially on the ray toward upper-left
infinity.  The same argument holds in the lower sector.

Here is an explicit factorial estimate for their common formal
series.  Set $x=y^{-2}=(2\tau)^{-1}$ and $D=x\partial_x$.
Equation \eqref{jump:q-equation} is formally equivalent to
\begin{equation}\label{jump:positive-recurrence}
 q=x+x\frac{2Dq-q^2}{(1-q)^2}
   =x+x\sum_{k\geq1}(2D-k+1)q^k.
\end{equation}
If $q=\sum_{n\geq1}b_nx^n$, then $b_1=1$ and
\begin{equation}\label{jump:coefficients}
 b_{n+1}=\sum_{k=1}^n(2n-k+1)
       \sum_{\substack{i_1+\cdots+i_k=n\\i_j\geq1}}
                       \prod_{j=1}^k b_{i_j}.
\end{equation}
Every summand is positive.  The $k=1$ term and induction give
the lower estimate in
\begin{equation}\label{jump:factorials}
 2^{n-1}(n-1)!\leq b_n\leq3^{n-1}(n-1)!.
\end{equation}
For the upper induction,
$\prod(i_j-1)!\leq(n-k)!$ and there are
$\binom{n-1}{k-1}$ positive compositions.  Hence
\[
 b_{n+1}\leq
 2\,3^{n-1}n!\sum_{j=0}^{n-1}\frac{3^{-j}}{j!}
 \leq 2e^{1/3}3^{n-1}n!<3^n n!,
\]
which proves \eqref{jump:factorials}.

We next verify the corresponding remainder estimate, rather than
infer it from coefficient growth.  Write $a_j=b_j/2^j$,
$P_N(\tau)=\sum_{j<N}a_j\tau^{-j}$, and
\[
 D_N=P_N'+P_N-2P_N^2+P_N^3-
                         \frac{1-2P_N}{2\tau}.
\]
As a polynomial in $\zeta=1/\tau$, $D_N$ vanishes to order $N$.
On $|\zeta|=(LN)^{-1}$, with $L$ fixed and large,
\eqref{jump:factorials} gives
$|P_N|+|\zeta^2\partial_\zeta P_N|\leq C/N$.
The maximum principle applied after division by $\zeta^N$
therefore gives
\[
 |D_N(\tau)|\leq (C/N)(LN)^N|\tau|^{-N}
                  \leq C A^N N!|\tau|^{-N}
 \quad (|\tau|\geq LN).
\]
The exact error $e_N=q_+-P_N$ obeys
\[
 e_N'+(1-\mathcal B)e_N=-D_N,
 \quad
 \mathcal B=2(q_++P_N)-(q_+^2+q_+P_N+P_N^2)-1/\tau.
\]
For $|\tau|\geq L'N$, choose $L'$ so large that along the
ray in \eqref{jump:volterra} one has
$\mathcal B=O(|\tau+tv|^{-1})$.  The integrating kernel is
bounded by $e^{-ct}(1+t/|\tau|)^C$.
Also
$|\tau|^N/|\tau+tv|^N\leq
 \exp(Nt/(c|\tau|))$.
The latter is absorbed by the exponential when $L'$ is large.
Variation from infinity proves
\begin{equation}\label{jump:factorial-remainder}
 |q_\pm(\tau)-P_N(\tau)|
       \leq C A^N N!|\tau|^{-N},\qquad |\tau|\geq L'N,
\end{equation}
uniformly on every closed smaller sector.  The boundary term at
infinity vanishes because the error is $O(|\tau|^{-1})$ there.

Suppose the lateral solutions agreed.  They would then give a
single holomorphic $q$ on an exterior sector of opening greater
than $\pi$, with \eqref{jump:factorial-remainder}.  Choose
$\pi/2<\omega<\vartheta$ and let $\Gamma_R$ run inward from
infinity on angle $-\omega$, counterclockwise along the right
circle $|\tau|=R$ to angle $\omega$, then outward to infinity.
For $|\arg u|<\omega-\pi/2$ set
\begin{equation}\label{jump:inverse-laplace}
 H(u)=\frac1{2\pi i}\int_{\Gamma_R} e^{\tau u}q(\tau)\,d\tau.
\end{equation}
The integral converges, is independent of sufficiently large $R$,
and is holomorphic near every positive $u$.  Move $R$ to $KN$,
with $K$ larger than the threshold in
\eqref{jump:factorial-remainder}.  Each monomial integral, closed
through the left half-plane, is its residue
$u^{j-1}/(j-1)!$.  The remainder on the arc and rays is at most
\[
 C A^N N!(KN)^{1-N}e^{KNu}
       \leq CKN(A/K)^Ne^{KNu}
\]
for positive $u$.  First take $K$ large, then $u>0$ small.
This tends to zero.  Consequently \eqref{jump:inverse-laplace}
equals the Borel series
\begin{equation}\label{jump:borel}
 \widehat q(u)=\sum_{n\geq1}
                         \frac{a_n}{(n-1)!}u^{n-1}
\end{equation}
near the positive origin and continues it analytically along the
entire positive axis.

This is impossible by the positive-coefficient singularity principle
known as Pringsheim's theorem
\cite[Theorem~IV.6 and Note~IV.13]{FS2009}, whose short argument
we give here.
The coefficients in \eqref{jump:borel}
are positive, and \eqref{jump:factorials} places its convergence
radius $R_H$ in $[2/3,1]$.  If it extended holomorphically through
$R_H$, choose a positive $r<R_H$ sufficiently close that its
Taylor disk extends to a positive $b>R_H$.  All Taylor
coefficients at $r$ are nonnegative, being termwise derivatives
of \eqref{jump:borel}.  Exchanging the resulting nonnegative
double sums at $b$ would show that the original series converges
at $b$, contradicting the definition of $R_H$.  Thus the two
lateral solutions are distinct.

Finally their slope difference satisfies the exact linear equation
\begin{equation}\label{jump:slope-difference-equation}
 (l_+-l_-)'=
 \left(l_++l_--\frac1y-
       \frac{l_+^2+l_+l_-+l_-^2}{y}\right)(l_+-l_-).
\end{equation}
The common expansion is
$l=y-y^{-1}-2y^{-3}+O(y^{-5})$, so the coefficient is
$-y+3/y+O(y^{-3})$.  Its integrable remainder gives
\eqref{jump:stokes-difference}; the constant is nonzero because
a nonzero solution of this scalar linear equation never vanishes.

For a finite bypass compare its slope with the canonical lateral
slope at the exit.  Their difference there is at most polynomial
in $Y$.  Along a path of length $O(Y)$ with $|y|\asymp Y$, the
coefficient in their exact difference equation is
$-y+o(Y)$, by \eqref{jump:translator-matching}.  Integrating
back to $y=Y$ therefore gives the factor
\[
 \exp\left(-Y^2/2+\Re y_{\rm e}^{2}/2+o(Y^2)\right).
\]
By \eqref{jump:negative-exit} this is bounded by
$\exp(-Y^2/2-c'Y^2)$.  Polynomial factors are absorbed by
decreasing $c'$.  This proves the final assertion and applies to
every subsequential exterior limit.
\end{proof}

\subsection{Boundary momentum and the common signed jump}

For a radial limit let
\begin{equation}\label{jump:lagrangian}
 \mathcal L(y,h,l)=y e^h\sqrt{1+l^2},
 \qquad \mathcal P(y,h,l)=\frac{ye^h l}{\sqrt{1+l^2}}.
\end{equation}
The square root is comparable to $y$ on the exterior arcs.
Direct differentiation using \eqref{jump:translator} gives
\begin{equation}\label{jump:primitive}
 \frac{d\mathcal P}{dy}=\mathcal L,
 \qquad
 d\mathcal P-\mathcal P\,dh
       =ye^h(1+l^2)^{-3/2}\,dl\quad(y\text{ fixed}).
\end{equation}
The second identity is the boundary-momentum cancellation.

At finite $Z_*$ the extra exponential in
\eqref{jump:physical-action} must be retained.  Indeed, with
$\varepsilon=Z_*^{-2}$, the radial integrand and its momentum are
\[
 \mathcal L_\varepsilon
 =ye^{h-\varepsilon(h^2+y^2)}\sqrt{1+l^2},\qquad
 \mathcal P_\varepsilon
 =ye^{h-\varepsilon(h^2+y^2)}\frac l{\sqrt{1+l^2}}.
\]
Their Euler--Lagrange identity is
$\frac{d}{dy}\mathcal P_\varepsilon
 =(1-2\varepsilon h)\mathcal L_\varepsilon$ along a stationary
radial curve; $\mathcal P_\varepsilon$ is not an exact primitive
of $\mathcal L_\varepsilon$.
For a noncircular finite-scale graph the boundary momentum instead
has denominator $\sqrt{1+h_y^2+h_\theta^2/y^2}$ and numerator
$ye^{h-\varepsilon(h^2+y^2)}h_y$, averaged in $\theta$ in the
units of \eqref{jump:physical-action}.
Thus we first use the exact finite-scale first variation, including
its height momentum.  Stationarity reduces the inner comparison to
its boundary at $y=Y$.  On that boundary and on bounded-$y$
exterior pieces, the coefficient correction is $o_B(1)$ uniformly
along the normalized homotopy, and the oscillation tends to zero.
Only there do we pass to the translator identities
\eqref{jump:primitive}.  The same preliminary matching value $Z$
and its derived scale $Z_*$ are held fixed along the homotopy;
the boundary height $h(Y)$ is allowed to vary, and its exact
momentum term is retained.  No polynomial bound for a gluing
graph is being compared directly with the exponentially small
action scale.

The outward surface orientation follows decreasing $r$, so its
positive-area integral is from the exit toward the incoming
section.  Thus a limiting exterior piece from $Y$ to its complex
exit contributes
$\mathcal P(Y)-\mathcal P(y_{\rm e})$ in the units preceding
the integrand in \eqref{jump:physical-action}.  The exit term
is $O(e^{-cY^2})$ by \eqref{jump:translator-matching} and
\eqref{jump:negative-exit}.

The inner action must be varied at the same time.  At finite $B$
truncate both normalized surfaces at their common $y=Y$ sections
and interpolate their outgoing derivative data there, on every
adverse end simultaneously.  Keep $p$ fixed.
Proposition~\ref{jump:compatibility} gives an actual normalized homotopy
whose endpoints are exactly the two restrictions.  At a fixed
section $r=2Y/Z_*$,
\[
 g_t=\frac{rl}{Z},
\]
so linear interpolation of the derivative data makes $l$ linear
in the interpolation parameter in the radial limit.  On the
nonadverse ends the paths, sections and derivative prescriptions are
fixed, while their boundary heights may vary.  Retain their fixed-exit
momentum terms in \cref{arr:first-variation}; their size is
$C_D B^M e^{-DB^2/5}$ by the same polynomial normalization bounds and
the full-stop Gaussian threshold.  This is negligible relative to
every adverse scale.  Stationarity and the exact first variation then
leave the adverse boundary momentum: in the units of
\eqref{jump:physical-action} its limit is $-\int\mathcal P\,dh$.
There is no compact multiplier term
because $dp=0$.  Combining this with the two exterior pieces,
\eqref{jump:primitive} gives the net comparison
\begin{equation}\label{jump:net-comparison}
 \int_{l_-(Y)}^{l_+(Y)}
                Ye^h(1+l^2)^{-3/2}\,dl+O(e^{-cY^2}).
\end{equation}
This calculation is made after $B\to\infty$ for fixed $Y$.
It therefore does not subtract two uncontrolled large limiting
pieces.  Smooth compact convergence and
\eqref{jump:interpolation-height}
justify first variation before this limit and passage to the
limit afterward.

Uniformly along the interpolation,
$h=Y^2/2-\log Y+o(1)$ and $l/Y=1+o(1)$ as $Y\to\infty$.
Consequently
\[
 Ye^h(1+l^2)^{-3/2}
       =Y^{-3}e^{Y^2/2}(1+o(1)).
\]
Lemma~\ref{jump:stokes} now shows that
\eqref{jump:net-comparison} tends to $c_0\ne0$.


\begin{proof}[Proof of Proposition~\ref{jump:jump}]
Lemma~\ref{jump:endpoints}, including the preliminary synchronization
\eqref{jump:nonadverse-sync} and the retained far boundary momenta,
changes the original actions to the
ones used in \eqref{jump:net-comparison}, at a cost
$O(e^{-cY_0^2})$ in the sum of absolute end scales.  At fixed
$Y_0$, \eqref{jump:physical-action} and the normalized homotopy
give the sum of the limiting comparisons, with errors $o_B(1)$
in each scale.  Lemma~\ref{jump:stokes} and
\eqref{jump:net-comparison} identify the same constant on every
end as $Y_0\to\infty$.  There are only finitely many ends, so
these errors sum with respect to their absolute scales.  The
order of limits is therefore the one stated.

There is no end-dependent orientation sign.  On every end $z$
increases outward along its chosen ray, the circular radius
decreases toward the collapsing profile, and the branch of area
starts positive.  The change to \eqref{jump:physical-action}
therefore always reverses the radial integration limits in the
same way.  Reversing a normal, or rotating an end's axis in
ambient space, changes none of these area conventions.  With
one common ordering of upper and lower bypasses, the same
$\epsilon_{\rm lab}c_0$ occurs throughout.  This proves the proposition.  The
separate phase comparison of the $Z_{*,j}$ will determine which
terms can dominate the sum; no assertion of cancellation or
noncancellation for arbitrary complex phases is built into
\eqref{jump:action-jump}.
\end{proof}
